\scanpage{17}
on a les isom fonctoriels
\[
u^*(F'\otimes G')\simeq u^*(F')\otimes u^*(G')
\]
et pour $u_*$ on a la ``formule de projection''
\[
u_*(F\otimes u^*(G'))\simeq u_*(F)\otimes G'
\]
[On a aussi des opérations $\Lambda^i$ (et autres opérations tensorielles) \unsure{que nous allons déduire} plus bas des opérations $\otimes$]

``Produit tensoriel externe''
\[
F\boxtimes G
\]
\margin{$F\in\operatorname{Ob}\mathcal A(\mathfrak X,G),\ G\in\operatorname{Ob}\mathcal A(\mathfrak X,H)$
\[
F\boxtimes G=\operatorname{pr}_1^*(F)\otimes\operatorname{pr}_2^*(G)\in\operatorname{Ob}\mathcal A(\mathfrak X,H\times G)
\]}
Est compatible avec les images inverses
\[
(u\times v)^*(F\boxtimes G)\simeq u^*(F)\boxtimes v^*(G)
\]
et redonne $F\otimes G$ par
\[
F\otimes G=(\operatorname{diag}_G)^*(F\boxtimes G)
\]
On pose
\[
\mathcal A(\mathfrak X,p)=\mathcal A(X,\unsure{\mathfrak S}_p)\qquad p\in\mathbb N
\]
\margin{\unsure{Éventuellement} introduire aussi $\mathcal A(\mathfrak X,p_1,\ldots,p_r)=\mathcal A(\mathfrak X,\mathfrak S_{p_1}\times\cdots\times\mathfrak S_{p_r})$}
Alors tout hom $u\deleted{\unsure{f_i}}:\mathfrak S_p\longrightarrow\mathfrak S_q$ donne des foncteurs
\diagram{\begin{tikzcd}
\mathcal A(\mathfrak X,p)&\mathcal A(\mathfrak X,q)\arrow[l,"u^*"']\\
\mathcal A(\mathfrak X,p)\arrow[r,"u_*"]&\deleted{\unsure{\mathcal E}}\mathcal A(\mathfrak X,q)
\end{tikzcd}}
\inserted{avec} \uline{transitivité} pour $u$
\margin{16}
